\section{Method}
\label{sec:method}

\subsection{Overview}
\label{sec:method:overview}

kam processes paired-end FASTQ files through six stages.
Stage~1 (\textbf{Parse}) extracts UMIs, skip bases, and template sequences from raw reads.
Stage~2 (\textbf{Assemble}) groups reads into molecule families by canonical UMI pair, detects UMI collisions via endpoint fingerprinting, and calls single-strand and duplex consensus sequences.
Stage~3 (\textbf{Index}) encodes consensus sequences as 2-bit k-mers and builds a molecule-level k-mer index.
Stage~4 (\textbf{Pathfind}) constructs a de~Bruijn graph from all on-target k-mers and enumerates paths between reference anchor k-mers using iterative depth-first search.
Stage~5 (\textbf{Call}) scores each path using molecule evidence, estimates variant allele frequency (VAF) with credible intervals, and applies strand bias and confidence filters.
Stage~6 (\textbf{Output}) emits called variants in VCF, TSV, and JSON formats with full provenance.

All stages operate on molecules as the atomic unit of evidence.
A k-mer supported by 100 PCR duplicates of one molecule contributes $n_\text{molecules} = 1$, not 100.
This distinction is the foundation of kam's error model.

% Figure~\ref{fig:pipeline} shows a horizontal flow diagram:
% FASTQ → Parse → Assemble → Index → Pathfind → Call → VCF
% with data types annotated on each arrow (ReadPair, Molecule, KmerIndex,
% GraphPath, ScoredPath, VariantCall).

\subsection{Read Parsing and UMI Extraction}
\label{sec:method:parsing}

kam targets Twist Biosciences duplex UMI chemistry~\autocite{Twist22cfdna}.
Both R1 and R2 follow read structure \texttt{5M2S+T} in fgbio notation~\autocite{fgbio}:

\begin{itemize}
  \item Positions 0--4: 5\,bp molecular barcode (UMI).
  \item Positions 5--6: 2\,bp skip/spacer (monotemplate, not biological).
  \item Position 7 onward: template (genomic) sequence.
\end{itemize}

Let $U_\text{R1}$ and $U_\text{R2}$ denote the 5\,bp UMIs extracted from R1 and R2.
The canonical UMI pair is:

\begin{equation}
  \label{eq:canonical-umi}
  (U_a, U_b) = \begin{cases}
    (U_\text{R1}, U_\text{R2}) & \text{if } U_\text{R1} \le U_\text{R2} \text{ lexicographically,} \\
    (U_\text{R2}, U_\text{R1}) & \text{otherwise.}
  \end{cases}
\end{equation}

Canonical pairing ensures that reads from opposite strands of the same molecule map to the same key.
In duplex sequencing, a forward-strand read pair carries $(U_\text{R1}, U_\text{R2}) = (\texttt{ACGTA}, \texttt{TGCAT})$, while the reverse-strand read pair from the same molecule carries $(\texttt{TGCAT}, \texttt{ACGTA})$.
Both produce the identical canonical pair $(\texttt{ACGTA}, \texttt{TGCAT})$.

If R1's UMI equals $U_a$, the read is classified as forward-strand. Otherwise it is reverse-strand. This strand assignment is used throughout consensus calling and variant scoring.

The 2\,bp skip bases serve as a quality-control signal.
Because the skip sequence is fixed within an adapter lot, reads whose skip bases deviate from the expected sequence are flagged.

\subsection{Molecule Assembly}
\label{sec:method:assembly}

Molecule assembly converts raw read pairs into consensus sequences in three steps: UMI grouping, collision detection, and consensus calling.

\subsubsection{UMI Grouping}

Reads are grouped by their canonical UMI pair (Equation~\ref{eq:canonical-umi}).
To tolerate sequencing errors in the UMI itself, kam applies directional Hamming clustering with distance threshold $d = 1$~\autocite{Smith17umitools}.
UMI groups within Hamming distance 1 are merged when one group's count is at least twice the other's, following the directional adjacency method of UMI-tools.

With 5\,bp random UMIs, there are $4^5 = 1{,}024$ possible values.
This limited diversity makes UMI collisions likely at high sequencing depths.

\subsubsection{Endpoint Fingerprinting}

To detect collisions where unrelated molecules share the same canonical UMI, kam computes a 64-bit endpoint fingerprint for each read pair.
The fingerprint encodes the first and last 8 bases of the template from both R1 and R2 using 2-bit encoding ($\texttt{A}=00$, $\texttt{C}=01$, $\texttt{G}=10$, $\texttt{T}=11$):

\begin{equation}
  \label{eq:fingerprint}
  F = \text{R1}_\text{first\,8} \| \text{R1}_\text{last\,8} \| \text{R2}_\text{first\,8} \| \text{R2}_\text{last\,8}
\end{equation}

\noindent where $\|$ denotes concatenation of 2-bit encoded subsequences.
Four endpoints at 8 bases each, 2 bits per base, fills exactly 64 bits, requiring no hashing.
Comparison is a single integer equality or XOR-plus-popcount operation.

Read pairs within the same UMI group whose fingerprints differ by more than 4 bits (approximately 2 base sequencing errors across all endpoints) are split into separate molecule families.
This threshold tolerates realistic error rates ($< 1\%$ family splitting at Q20) while catching genuine collisions.

\subsubsection{Single-Strand Consensus}
\label{sec:method:ssc}

Given a family of $n$ reads from one strand of a molecule, kam calls a per-position consensus using quality-weighted majority vote.

At position $i$, let $Q_{i,j}$ be the Phred quality of the $j$-th read.
The error probability is:

\begin{equation}
  \label{eq:phred-to-prob}
  p_{i,j} = 10^{-Q_{i,j}/10}
\end{equation}

The weight assigned to read $j$'s base call at position $i$ is $w_{i,j} = 1 - p_{i,j}$.
Bases with $Q_{i,j}$ below a minimum threshold (default: $Q = 10$) are excluded.
For each candidate base $b \in \{A, C, G, T\}$, the total weight is:

\begin{equation}
  \label{eq:weight-sum}
  W_b(i) = \sum_{\substack{j=1 \\ \text{base}_{i,j} = b}}^{n} w_{i,j}
\end{equation}

The consensus base at position $i$ is:

\begin{equation}
  \label{eq:consensus-base}
  \hat{b}(i) = \operatorname*{argmax}_{b} \; W_b(i)
\end{equation}

The estimated error probability of the consensus at position $i$ is:

\begin{equation}
  \label{eq:consensus-error}
  \hat{p}_\text{SSC}(i) = 1 - \frac{W_{\hat{b}(i)}(i)}{\sum_{b} W_b(i)}
\end{equation}

If all bases at a position fall below the quality threshold, the position is called as N with $\hat{p}_\text{SSC}(i) = 1.0$.
Output consensus quality is capped at a configurable maximum (default: Q60).

\subsubsection{Duplex Consensus}
\label{sec:method:duplex}

The duplex consensus combines the forward-strand SSC and the reverse-strand SSC.
At each position $i$, let $\hat{b}_\text{fwd}(i)$ and $\hat{b}_\text{rev}(i)$ be the single-strand consensus bases, with error probabilities $\hat{p}_\text{fwd}(i)$ and $\hat{p}_\text{rev}(i)$.

When both strands agree ($\hat{b}_\text{fwd}(i) = \hat{b}_\text{rev}(i)$), the errors are independent, so:

\begin{equation}
  \label{eq:duplex-error}
  \hat{p}_\text{duplex}(i) = \hat{p}_\text{fwd}(i) \times \hat{p}_\text{rev}(i)
\end{equation}

This multiplication yields error rates near $10^{-7}$ for typical input qualities~\autocite{Schmitt12duplex}, which is what makes detection at 0.1\% VAF feasible.

When the strands disagree ($\hat{b}_\text{fwd}(i) \neq \hat{b}_\text{rev}(i)$), the default behaviour picks the base from whichever SSC had lower error probability and marks the position as non-duplex.
A strict mode (\texttt{-{}-strict-duplex}) masks the position as N.
In either case, the position does not receive duplex-level confidence.

Each assembled molecule carries a \texttt{MoleculeEvidence} record that tracks the number of forward-strand reads, reverse-strand reads, and whether the molecule has duplex support.
This record propagates through all downstream stages.

\subsection{K-mer Indexing}
\label{sec:method:index}

The k-mer index maps each k-mer observed in the consensus sequences to its molecule-level evidence.

\subsubsection{2-Bit Encoding}

Each DNA base is encoded as 2 bits: $\texttt{A}=00$, $\texttt{C}=01$, $\texttt{G}=10$, $\texttt{T}=11$.
A k-mer of length $k$ occupies $2k$ bits in a \texttt{u64}, supporting $k \le 31$.
Bases are packed left-to-right, with the first base in the most significant bit pair.

\subsubsection{Canonical K-mers}

Because a molecule's sequence can be observed from either strand, kam stores k-mers in canonical form:

\begin{equation}
  \label{eq:canonical-kmer}
  \text{canonical}(s) = \min\bigl(s, \;\overline{s}^{\,R}\bigr)
\end{equation}

\noindent where $\overline{s}^{\,R}$ denotes the reverse complement of $s$.
In 2-bit encoding, the complement of each base is obtained by XOR with \texttt{0b11}, and the base order is reversed.

\subsubsection{Sliding Window}

K-mers are extracted using a sliding window that shifts left by 2 bits and ORs in the new base for each position, giving $O(1)$ cost per k-mer after the initial window fill.
When an ambiguous base (N) is encountered, the window resets and skips forward until $k$ consecutive valid bases are available.

\subsubsection{Molecule-Level Accumulation}

For each canonical k-mer, the index stores a \texttt{MoleculeEvidence} record:

\begin{itemize}
  \item $n_\text{molecules}$: number of distinct molecules carrying this k-mer.
  \item $n_\text{duplex}$: of those, how many have duplex support.
  \item $n_\text{simplex\_fwd}$, $n_\text{simplex\_rev}$: molecules with only forward or reverse strand support.
  \item $\bar{p}_\text{error}$: mean per-base error probability across supporting molecules.
\end{itemize}

Each unique molecule contributes 1 to $n_\text{molecules}$ regardless of how many PCR duplicates it produced.
In targeted mode, only k-mers matching an allowlist (derived from reference and expected variant sequences, padded by $k-1$ bases) are indexed.

\subsection{De Bruijn Graph Construction and Path Walking}
\label{sec:method:pathfind}

\subsubsection{Graph Construction}

kam builds a de~Bruijn graph~\autocite{Iqbal12cortex} from on-target k-mers.
Nodes are raw (non-canonical) k-mers, stored in the original read orientation so that consecutive k-mers maintain their suffix-prefix overlap.
Both the forward and reverse-complement of each on-target k-mer are included as nodes.
A directed edge from node $A$ to node $B$ exists when the $(k{-}1)$-base suffix of $A$ equals the $(k{-}1)$-base prefix of $B$:

\begin{equation}
  \label{eq:edge-condition}
  A[1{:}k{-}1] = B[0{:}k{-}2]
\end{equation}

In 2-bit encoding, this is:

\begin{equation}
  \label{eq:edge-2bit}
  (A \;\texttt{AND}\; M_{k-1}) = (B \gg 2)
\end{equation}

\noindent where $M_{k-1} = 2^{2(k-1)} - 1$ is the bitmask for the lower $2(k{-}1)$ bits and $\gg$ denotes right shift.

The graph is built from all k-mers observed in on-target molecules, not only reference k-mers.
Variant k-mers (those differing from the reference by one or more bases) must be nodes in the graph for their paths to be discoverable.

\subsubsection{Anchor K-mers}

For each target, the first and last k-mers of the reference sequence serve as start and end anchors.
All paths are enumerated from the start anchor to the end anchor.
The reference path passes through all k-mers of the wildtype sequence.
Variant paths diverge from the reference at the site of the mutation.

\subsubsection{Path Walking}

Algorithm~\ref{alg:dfs} describes the DFS path enumeration.

\begin{algorithm}
\caption{Iterative DFS path enumeration between anchor k-mers.}
\label{alg:dfs}
\begin{algorithmic}[1]
\Require{Graph $G$, start anchor $s$, end anchor $t$, max path length $L$, max paths $P$}
\Ensure{Set of paths $\mathcal{P}$ from $s$ to $t$}
\State $\mathcal{P} \gets \emptyset$
\State $\mathit{stack} \gets [(s, 0)]$ \Comment{(node, next-successor-index)}
\State $\mathit{path} \gets [s]$
\State $\mathit{visited} \gets \{s\}$
\While{$\mathit{stack}$ is not empty \textbf{and} $|\mathcal{P}| < P$}
  \State $(v, i) \gets \mathit{stack}.\text{last}()$
  \If{$i \ge |\text{successors}(v)|$} \Comment{All successors tried: backtrack}
    \State $\mathit{stack}.\text{pop}()$; $\mathit{path}.\text{pop}()$; $\mathit{visited}.\text{remove}(v)$
    \State \textbf{continue}
  \EndIf
  \State $u \gets \text{successors}(v)[i]$; increment $i$ in stack
  \If{$u \in \mathit{visited}$ \textbf{or} $|\mathit{path}| \ge L$}
    \State \textbf{continue} \Comment{Cycle or path too long}
  \EndIf
  \If{$u = t$}
    \State $\mathcal{P} \gets \mathcal{P} \cup \{\mathit{path} \cup \{u\}\}$
  \Else
    \State $\mathit{visited}.\text{insert}(u)$; $\mathit{path}.\text{push}(u)$; $\mathit{stack}.\text{push}((u, 0))$
  \EndIf
\EndWhile
\State \Return $\mathcal{P}$
\end{algorithmic}
\end{algorithm}

The DFS uses a single path buffer and a \texttt{HashSet} for O(1) cycle detection.
Memory use is $O(B \times L)$ where $B$ is the maximum branching factor and $L$ is \texttt{max\_path\_length}, constant regardless of the number of paths enumerated.

In contrast, a BFS implementation clones the full partial path at every branch, accumulating $O(B^{L/2} \times L)$ bytes in the queue at peak.
At 1 million reads and 30\,ng input, the BFS queue saturates available memory within the pathfind stage for several targets in the titration panel.
The DFS implementation eliminates this scaling problem entirely.

The maximum path length is set per target:
\begin{equation}
  \label{eq:max-path-length}
  L = (T - k + 1) + 50
\end{equation}
\noindent where $T$ is the target length in bases.
For a 100\,bp target with $k = 31$, this gives $L = 120$: the 70-k-mer reference path plus 50 k-mers of headroom for large indels.
The constant default ($L = 500$) in earlier versions allowed the walker to explore paths far outside the target window, greatly amplifying the BFS memory cost.

Sequences are reconstructed from paths by taking all $k$ bases from the first k-mer and appending the last base of each subsequent k-mer, yielding a sequence of length $k + |\pi| - 1$ from a path of $|\pi|$ k-mers.

\subsection{Statistical Variant Calling}
\label{sec:method:calling}

\subsubsection{Path Scoring}

Each path is scored by looking up the \texttt{MoleculeEvidence} for every constituent k-mer.
The aggregate statistics are:

\begin{itemize}
  \item $n_\text{min}$: minimum $n_\text{molecules}$ across all k-mers in the path.
  \item $\bar{n}$: mean $n_\text{molecules}$ across all k-mers in the path.
  \item $d_\text{min}$: minimum $n_\text{duplex}$ across all k-mers.
  \item $s_\text{fwd}$, $s_\text{rev}$: minimum forward and reverse simplex molecule counts across all k-mers in the path.
\end{itemize}

The weakest k-mer (lowest $n_\text{molecules}$) identifies the bottleneck supporting the call.
The path whose reconstructed sequence matches the reference is marked as the reference path.

\subsubsection{VAF Estimation}

Let $n_a$ be the number of molecules supporting the alternate path (from $n_\text{min}$ of the alt path) and $M$ be the total molecules ($n_\text{min,ref} + n_\text{min,alt}$).
The point estimate of VAF is:

\begin{equation}
  \label{eq:vaf}
  \hat{f} = \frac{n_a}{M}
\end{equation}

The 95\% credible interval uses the conjugate Beta posterior with a uniform prior $\text{Beta}(1, 1)$.
Observing $n_a$ successes in $M$ trials yields posterior $\text{Beta}(n_a + 1, M - n_a + 1)$:

\begin{equation}
  \label{eq:vaf-ci}
  \text{CI}_{95\%} = \bigl[q_{0.025}, \; q_{0.975}\bigr]
  \quad \text{where} \quad
  q_\alpha = F^{-1}_{\text{Beta}(n_a+1, \, M-n_a+1)}(\alpha)
\end{equation}

\noindent and $F^{-1}$ is the inverse CDF (quantile function) of the Beta distribution.

\subsubsection{Strand Bias}

kam tests for strand bias using Fisher's exact test on a $2 \times 2$ contingency table.
Let $a = s_\text{fwd,alt}$, $b = s_\text{rev,alt}$ be the minimum forward and reverse simplex molecule counts for the alternate path, and $c = s_\text{fwd,ref}$, $d = s_\text{rev,ref}$ the corresponding counts for the reference path.
These are per-molecule counts (the path minimum across k-mers), not per-k-mer sums:

\begin{equation}
  \label{eq:strand-bias-table}
  \begin{array}{c|cc}
    & \text{Forward} & \text{Reverse} \\ \hline
    \text{Alt}  & a & b \\
    \text{Ref}  & c & d
  \end{array}
\end{equation}

The two-tailed p-value is computed by summing all hypergeometric probability masses $\le$ the observed table's probability, given fixed marginals.
Variants with $p < 0.01$ are flagged as strand-biased.

\subsubsection{Confidence Score}

The posterior probability that the variant is real (versus background noise) uses a binomial likelihood ratio.
Under equal priors, the posterior is:

\begin{equation}
  \label{eq:confidence}
  P(\text{signal} \mid \text{data}) = \frac{L(\hat{f})}{L(\hat{f}) + L(\varepsilon)}
\end{equation}

\noindent where $L(p) = p^{n_a} (1 - p)^{M - n_a}$ is the binomial likelihood and $\varepsilon$ is the per-site background error rate (default: $10^{-4}$).
Variants require a posterior probability $\ge 0.99$ to pass the confidence filter.

\subsubsection{Variant Classification}

Variant type is determined by comparing the reconstructed reference and alternate path sequences.
Let $|R|$ and $|A|$ denote their lengths:

\begin{itemize}
  \item $|R| = |A|$ with exactly 1 differing base: SNV.
  \item $|R| = |A|$ with $>1$ differing bases: MNV.
  \item $|R| > |A|$: deletion.
  \item $|R| < |A|$: insertion.
\end{itemize}

\subsubsection{Filter Assignment}

Each variant receives one of five filter labels.
A variant passes all filters when: the alt molecule count $n_a \ge 2$ (configurable), the duplex alt molecule count $\ge 1$ (configurable), the strand bias p-value $\ge 0.01$, and the posterior confidence $\ge 0.99$.
Filters are evaluated in this order; a variant failing any criterion is labelled with the first filter it fails: \texttt{LowConfidence}, \texttt{LowDuplex}, \texttt{StrandBias}, or \texttt{CollisionRisk}.

Requiring at least one duplex alt molecule suppresses false positives from background sequencing errors.
A random single-base error producing 2+ simplex alt molecules is unlikely to produce independent errors on both strands of the same molecule, so it cannot generate duplex evidence.
True variants at $\ge 0.1\%$ VAF with 5--15\% duplex rates will have duplex alt evidence in all but the rarest cases.

\subsubsection{VCF Allele Extraction and Left-Normalisation}

The graph path walk produces full 100\,bp reference and alternate sequences.
Converting these to standard VCF alleles requires three steps.

First, the minimal allele is extracted: the common prefix and suffix shared by the reference and alternate sequences are trimmed, leaving only the differing region.
For SNVs and MNVs (sequences of equal length) this is the final representation.
For indels, the inner common prefix and suffix of the differing region are further stripped to isolate the inserted or deleted bases.

Second, indels are left-normalised to match the standard VCF convention used by GATK, bcftools, and reference truth sets~\cite{tan2015}.
An anchor base is placed immediately before the event.
While the preceding reference base matches the last base of the event sequence, the event sequence is rotated right by one base and the anchor shifts one position to the left.
This places the variant at the leftmost possible position in any repeat context.

Third, the anchor base is prepended as required by the VCF specification: for a deletion, REF\,=\,anchor\,+\,deleted\,bases and ALT\,=\,anchor; for an insertion, REF\,=\,anchor and ALT\,=\,anchor\,+\,inserted\,bases.

Genomic coordinates are derived from the target header (\texttt{chrN:START-END}), where START is the 1-based genomic position of the first base in the stored sequence. A sequence index $i$ (0-based) maps to VCF position $\text{START} + i$.
